\chapter{The Residue Refuses To Disappear}

\section{The antagonist and the treasure}

The last chapter's loop kept the fact --- but keeping it was not free. To keep the fact, the round-trip had to record
how far the fact had to travel to be kept, and that distance does not vanish when the trip is over. It sits there, at
the bottom of every refinement, the one thing the climb could not climb away. This chapter names it. The name is the
\emph{residue}, and it is two things at once, which is why it takes a chapter to introduce. It is the
\emph{antagonist} --- the term no refinement can discharge, the leftover the machine would frankly rather not have to
measure --- and it is the \emph{treasure}, because the leftover \emph{is} the reading. Everything the instrument is
built to measure, it will read off exactly this.

The residue is not a new primitive smuggled in from outside; it is a type at the top of the tower the book has been
building. In this construction the residue's type \emph{demands} every earlier tier before it will form at all --- the
distinguishable, the admissible, the countable, the encoded, each required as a standing condition --- and only once
all four are in hand does it add anything of its own. So the residue is, definitionally, what remains once all four
refinements have run: the tier you reach by refining, and can refine no further. ``Un-eliminable'' is a claim about
the \emph{type}, not a mood --- it means there is no tier above this one to refine it into. The ladder built so far
ends on the leftover.

A word on ``so far,'' because the book means to keep its promises exact. The residue is the top of \emph{this} tower
--- the one the first three chapters climb --- and not the top of everything. Further on, the construction will stack
more tiers above the residue and carry it upward, and the next chapter's seam is precisely that carrying. So read
``un-eliminable, the top'' as top-\emph{so-far}, the summit of the climb to here, and not as a final ceiling. The
machine that, in its own words, does not much care to measure leftovers discovers that the leftover is the summit of
its own construction.

What the residue's type adds, once it has demanded the tower beneath it, is two fields, and the second is the one to
watch. It carries a relation --- a predicate on the settled limits, saying when one settled thing stands to another as
the residue relates them --- and the shape of that relation is not new. It is the loop of the last chapter, lifted one
floor up. Where the round-trip compared raw sequences, this compares settled limits; but the cases are the same. At
the bottom, the two facts must be equal. While both are still refining, the fact is held equal, a scale is permitted
to rise, and an index is permitted to deepen. The residue, in other words, is measured in exactly the fact-preserving
currency the loop established: the fact conserved, a magnitude climbing, a direction deepening. And here the seam the
last chapter planted finally pays out. An encoding, that chapter observed, hands you an origin and a limit away from
it --- a \emph{direction} and a \emph{magnitude}, one covariant and one contravariant. The residue is that pair, made
into a type: the conserved fact and its index are the direction, the rising scale is the magnitude, and the two
together are the whole of what the leftover carries.

Now the turn the chapter is built on: the antagonist is the treasure. Why climb a whole tower of refinements only to
be left, at the top, with a term you cannot discharge? Because that term is the measurement. It is not the failure of
the construction; it is the point of it. Everything downstream reads this one object and nothing else: a later chapter
will bracket it, another will name the box it falls into, another will put the coupling-question to it, and the book
will end by holding it open as the honest interval it can only bracket. The residue is what the whole instrument
converges on. And the thing about a treasure the machine would rather not measure is that it will not leave on its own
--- the next section shows that discarding it is not even a permitted operation, that the attempt to drop the leftover
fails to typecheck. What refuses to disappear is exactly what the instrument will spend the rest of the book reading.

It is worth seeing the residue once in the two verbs, because it is where they finally point. Bag the differences by
their admissibility, lay the bag out in an enumeration, encode it, let it settle --- funge upon funge upon funge ---
and there is one characteristic that will not bag away, that survives every gathering: the residue. It is the tange no
funge removes. It was selected, at the very start, by the whole tower at once, and no amount of bagging by likeness
can dissolve it back into nothing; and that irreducibility is precisely why it is the treasure. The count could be
fungeed; the limit could be fungeed; the leftover cannot. (That the residue \emph{is} a direction and a magnitude, an
origin and a distance from it, is the reading the book lays over the relation --- and the source, tellingly, flags the
reach of its own framing here more loudly than anywhere before it, which is exactly the signal to hold the reading as
a reading. The thing the code builds is the tier and its relation; the direction-and-magnitude is what they are read
as.)

So the third chapter opens on its antagonist, which is its treasure, and the whole of what follows is the study of
what this one un-dischargeable term turns out to be. The book has climbed a tower in order to arrive, deliberately, at
the thing it could not climb past --- and that thing, it will argue for the rest of these pages, is quietly almost
everything the instrument can measure.

\section{The lie of discarding}

The section before this named the residue un-eliminable and left the word standing as a claim about a type: there is
no tier above it to refine it into. This section shows what the machine does when you reach past that refusal and try
to eliminate the residue anyway --- when you attempt to throw the leftover away. The answer is sharper than ``you lose
something.'' The answer is that the attempt is \emph{false}. Not discouraged, not merely lossy, not inadvisable:
false, in the exact sense that the proposition asserting it has no proof, and the checker that would have to accept
that proof cannot find one.

To see why, look at what ``discard the residue'' means in the machine's own terms. Every comparison the construction
makes is an order --- a relation that says when one term stands below another. Each such relation is defined by cases
on the two terms it compares, and among those cases is the one that decides what happens when a populated, indexed
term is asked to stand below \emph{nothing} --- below the empty term the construction writes \verb|.nil|. To discard
the residue is precisely to assert that step: that the indexed leftover reduces back down to the empty term, as though
it had never been drawn. And the relation's verdict on that case is a single word: \verb|False|.

That is not a figure of speech. The relation is a proposition, and a proposition is inhabited or it is not; ``this
populated term relates down to nothing'' is the uninhabited one. There is no proof of it, because the case that would
have to supply one returns \verb|False| outright. Where the loop of the previous chapter and the opening of this one
read the \emph{diagonal} cases --- the ones where the fact is held equal and the climb is permitted --- this section
reads the \emph{off-diagonal} case, the one asked when a made distinction is pointed back at nothing. The residue
refuses to disappear for the most literal reason a thing can refuse: the sentence ``it disappeared'' names no
inhabitant. There is nothing to check, so the checker declines.

Now set that case beside its partner. The same relations answer the opposite question --- what happens when
\emph{nothing} is asked to stand below a populated term --- and there the verdict inverts to \verb|True|. The empty
term stands below anything; the source, at the place it settles this, jokes that it is letting zero in as a converged
digit the way a famous predecessor once let something in, and then rules plainly that zero is always less than
anything, converged or not. So the two cases together carve a direction into the order. Upward --- from nothing toward
a made distinction --- is always admissible. Downward --- from a made distinction back toward nothing --- is always
\verb|False|. Refinement adds, and is always allowed; discarding subtracts, and is never allowed. That one-wayness is
not a decoration on the construction; it \emph{is} the construction's soundness. A machine that could descend that arm
could quietly un-make a distinction it had already drawn, and a reading that can be quietly un-made can be faked. The
honesty the whole book is climbing toward rests on this one asymmetric pair: you may build up from nothing; you may
never sink back to it.

And here is the part worth slowing for: this is not a rule about the residue. It is a rule the machine obeys
\emph{everywhere}. The same two cases --- \verb|False| going down, \verb|True| coming up --- stand in the order on raw
sequences, in the order on settled limits, in the round-trip encoding the loop chapter built, and in the residue's own
relation at the top. Four relations, from the bottom of the tower to its summit, and one discipline running through
all of them: nothing populated is ever less than nothing. So the chapter title --- the residue refuses to disappear
--- is, read exactly, the standing law of every comparison the machine makes. The antagonist of this chapter is not a
special leftover with a special stubbornness; it is one instance of a house rule the construction has enforced since
its first sequence. The residue simply happens to be the term the rule protects at the top. (A small honesty about the
foot of that ladder: the two order relations let the empty term sit at or below \emph{another} empty term for free,
while the two encodings ask, there, that the bare facts still match --- so it is the \emph{discard} case, and its
upward partner, that run uniform across all four, not every case alike. The claim is carried by the arm that carries
it.)

There is one more thing the code says out loud at exactly this rule, and because it says it in a comment I can hand it
over as derived rather than read in. On the diagonal case --- where both terms are populated --- the two orders, on
sequences and on limits, split into two branches, and the source annotates them by name. Where the facts \emph{agree},
the comparison runs one way and is marked \emph{covariant}; where the facts \emph{differ}, it runs the other way, the
order on the scale inverted, and is marked \emph{contravariant}. This is the book's two verbs written into the order
itself. To compare with the fact held equal --- to bag two terms by the likeness of what they carry --- is to funge,
and the source calls it covariant. To compare with the fact differing --- to keep two terms apart by the very
characteristic that distinguishes them --- is to tange, and the source calls it contravariant. The two bands that a
later chapter will open into four faces are already here, cut into the order rule this section reads. And the source's
confidence in the cut is worth noticing: the hedge it wrote large a section ago, when it read the residue as a
direction and a magnitude, drops to its quietest here, at the order rules, where the source is sure. The deeper the
reading, the louder the machine's own doubt; at the plain rules, it is quiet. That gradient is a tool --- it tells you
which of the book's claims sit on firm ground and which are the gauge reaching.

Which lets me say what discarding is in the two verbs. To drop the residue is to try to un-funge and un-tange in one
step: to pull the member out of the bag and, at the same instant, erase the selection that put it there --- to pretend
the distinction was never drawn. The \verb|False| arm is the machine declining exactly that pretense. You tanged it;
you do not now get to claim you didn't. The bag only grows; no member is ever un-selected back out of it. And a
caution on the word this section keeps circling: it is tempting to call the discard an ill-typed \emph{coercion}, a
thing the compiler forbids by construction --- but there is no coercion in the code, no such instance anywhere in the
source, only a proposition that returns \verb|False|. ``Coercion'' and ``lie'' are the reading; the thing built is the
uninhabited case. The book keeps the two apart, as always.

The forward sentence of this section is one the whole book will lean on at the end. The reason the jar can be honest
--- the reason the interval it finally holds open cannot be a crank's fabricated point --- is this \verb|False| arm. A
machine that cannot descend the discard direction cannot un-make a measurement, and a measurement that cannot be
un-made cannot be faked. The open bracket the last chapter will reach is, in the end, exactly this: the residue not
discarded --- held, rather than thrown away. Honesty, in this construction, has a precise type. It is a \verb|False|
arm: the impossibility of pretending you did not measure what you measured.

\section{What the residue becomes}

The last section closed one direction: the residue cannot be discarded --- the step down to nothing is \verb|False|,
and the machine will not compute it. So the machine does the only other thing it can. It cannot throw the leftover away
and it cannot store it whole, so it takes a \emph{bite} of the residue and hands the bite on. In the source's own
dialogue, the bite proves nothing --- but it gives the next machine something small enough to lie about. That bite has
a name, and it is the type this section is about: the \verb|Sample|. What the residue \emph{becomes}, once the machine
has refused both to keep it whole and to drop it, is a small finite datum, passable to whatever reads next.

It is worth saying at the outset that the source is most nervous exactly here. Each type in this tower carries a
private hedge --- the running self-flag the last section read as a fidelity gradient --- and the \verb|Sample|'s is the
highest anywhere in the construction: higher than the residue's, higher than the encoding's, far above the order rules
the source was sure of. The reason is plain once named. This is the type that most wants to be \emph{a number}, and the
source knows it, and says so by hedging loudest exactly where the reach is longest. Read that as instruction.
Everything in this section that sounds like ``the residue is a number now'' is the reading, and the machine's own doubt
is turned up to its maximum to tell you so. What is actually built is smaller and firmer, and the section keeps to it.

What is actually built is an inductive type with two constructors, and the two together say the residue becomes not a
value but a \emph{process that begins and then answers itself}. The first constructor is a seed: from a fact and a
settled limit it forms the \verb|Sample| at its origin --- the starting datum you cannot get out from under needing.
(The source nods here to Dirichlet, from far back: you cannot, it says, escape the need for some magical number to
begin with; the seed is that unavoidable beginning.) The second constructor is a response: given fresh facts and limits
and a \emph{prior} \verb|Sample|, it builds the next one atop it. So a \verb|Sample| is a seed and then a recursion ---
a thing that starts, and then responds to what it has already said. The residue, refused as a whole and refused as
nothing, comes back as a datum that continues itself.

The order on these samples is where the two verbs, for the first time, run \emph{at once}. Recall that the order rules
of the last section split into a covariant branch, where the facts agree, and a contravariant branch, where they
differ. The \verb|Sample| order re-runs that split --- but now it carries two facts and two scales, and routes them
separately. The source poses it as a question and then takes the greedy answer: covariant on the one, contravariant on
the other --- why not both? So a single comparison of two samples funges on the fact that agrees and tanges on the fact
that differs, in the same breath. This is the double-read the book has been pointing toward: one datum, counted
\emph{with} itself along the characteristic it shares and \emph{against} itself along the characteristic it does not.
The residue, once sampled, carries both bands in one order --- and that single object, read on two bands at once, is the
exact seed of the four faces a later chapter will open. The faces do not come from four different things; they come from
one sample read covariantly and contravariantly together.

And now the word the whole book will end on. Among the \verb|Sample| order's cases is one that must never hold: a
response standing below its own origin --- a continuation claiming to be less than the seed it grew from. The relation's
verdict on it is \verb|False|, the same refusal the last section studied; but here the source stops and \emph{names} the
refusal. That collapse, its comment says, should not happen without serious \emph{strain} --- the compiler thought
\emph{true = false} somewhere. And this is not a stray word chosen for color. A few lines on, the source formally
introduces the concept under its own heading --- \emph{informational strain} --- and glosses it precisely: informational
strain is \emph{true \(\neq\) false} forced to \emph{true = false}. So the forbidden collapse of this datum \emph{is}
informational strain: the attempt to make a distinction and its own negation hold together, the one thing a sound
checker computes to \verb|False|. This matters past the third chapter, because \emph{strain} is the word the last
chapter of the book turns on --- and it is not, it turns out, a metaphor the final gauge reached for. It is a concept the
construction introduces here, at the close of Part One, in its own voice. The jar the book ends by holding open is the
strain \emph{relaxed}: the honest machine, offered the collapse, declines to force \emph{true = false} and holds the
interval open instead. This section plants that through-line at its root --- the residue becomes a \verb|Sample| whose
single illegal step is exactly the strain the whole instrument, at the end, refuses to force.

One boundary before the chapter closes, and it is the null-basis boundary this series is built on. This sample is the
datum every later reading will interpret --- and the physical gauge, the next volume, will read it as mass, as charge,
as phase, as cost. Those four names are Volume Two's cut, not this one's. This volume claims only the type: the residue
becomes a \verb|Sample|, a seed with a recursive response, doubly routed on the two bands, forbidden the one collapse
the source calls strain. The physical readings are the interpretation the next machine --- in the source's wry sense ---
is handed something small enough to lie about. This volume hands the sample over and does not itself read it. The
residue has become a passable symbol; what it means is the next book's to say.

So the third chapter closes where it must. The antagonist arrived un-dischargeable; the refusal proved itself a
\verb|False| arm; and the leftover --- kept and passed rather than stored or dropped --- became a \verb|Sample|
carrying, in a single order, both bands and one forbidden collapse. The word for that collapse is the word the book
ends on. The residue is not yet read; it will not be read until far later, and even then it will come back, honestly, as
an interval. But it is now a thing that can be handed on --- which is all the machine, having refused to lose it, ever
needed it to be.
